% =====================================================================
% Kapitel 5: Vendor-Adapter
% =====================================================================

\chapter{Vendor-Adapter}

\label{ch:vendor}

\emph{Vendor-Adapter} sind die Komponenten, die mit den externen
Datenquellen sprechen. Jeder Adapter hat die gleiche Schnittstelle
-- \code{fetch(symbol, start, end)} und \code{healthcheck()} -- aber
unterschiedliche Implementierungen je nach API-Protokoll der
Datenquelle.

\section{Das Vendor-Interface}

\begin{definition}
\textbf{Vendor-Adapter}
Ein \emph{Vendor-Adapter} ist eine Python-Klasse, die eine bestimmte
externe Datenquelle (z.\,B.\ FRED) abstrahiert. Sie bietet eine
\emph{einheitliche} Schnittstelle, damit der Rest des Systems nicht
wissen muss, ob die Daten von einer REST-API, einer SDMX-Quelle
oder einer CSV-Datei kommen.
\end{definition}

Listing~\ref{lst:vendor-base} zeigt die
\emph{Abstrakte-Basisklasse}, von der alle konkreten Adapter erben.

\begin{lstlisting}[language=python, caption={Basisklasse aller Vendor-Adapter},label=lst:vendor-base]
from abc import ABC, abstractmethod
from datetime import date
import polars as pl

class Vendor(ABC):
    """Base class for all vendor adapters."""
    code: str              # e.g. "fred", "ecb", "yahoo"
    base_url: str          # Vendor's API base URL

    @abstractmethod
    def fetch(self, symbol: str, start: date, end: date) -> pl.DataFrame:
        """Fetch raw data for `symbol` in [start, end] inclusive."""
        raise NotImplementedError

    @abstractmethod
    def healthcheck(self) -> bool:
        """Return True if the vendor is reachable."""
        raise NotImplementedError
\end{lstlisting}

Die Wahl von \code{polars.DataFrame} als Rückgabetyp ist
bewusst: Polars ist schnell, hat Lazy-Evaluation und ist vollständig
typisiert. Die Adapter sind dünn -- sie übersetzen JSON oder XML in
ein DataFrame -- und die schwere Arbeit (Statistiken, Resampling,
Joins) passiert downstream in dedizierten Analytics-Modulen
(siehe Kapitel~\ref{ch:analytics}).

\section{FRED -- Federal Reserve Economic Data}

\begin{definition}
\textbf{FRED}
\emph{FRED} (Federal Reserve Economic Data) ist die
Wirtschaftsdatenbank der US-Notenbank Federal Reserve. Sie enthält
über 800.000~Zeitreihen zu Zinssätzen, Inflation, Beschäftigung,
BIP, Geldmenge und vielem mehr. Die Daten sind \foreign{public domain}
und über eine kostenlose REST-API zugänglich.
\end{definition}

\begin{definition}
\textbf{API-Key}
Ein \emph{API-Key} ist ein geheimer Token, der den
\emph{API-Aufrufer} identifiziert. FRED verlangt einen Key, den man
kostenlos unter \url{https://fred.stlouisfed.org/docs/api/api_key.html}
beantragen kann. Der Key wird in der Datei \code{.env} (nie im
Quellcode!) gespeichert.
\end{definition}

FRED liefert Zeitreihen über die Endpoints
\code{/fred/series/observations} (für eine einzelne Zeitreihe) und
\code{/fred/series} (für Metadaten). Listing~\ref{lst:fred-fetch}
zeigt, wie ein typischer Fetch aussieht.

\begin{lstlisting}[language=python, caption={FRED-Vendor: Healthcheck und Fetch},label=lst:fred-fetch]
import httpx
from deephold_db.utils.rate_limit import limit
from deephold_db.utils.retry import http_retry
from deephold_db.utils.config import get_settings

class FredVendor(Vendor):
    code = "fred"
    base_url = "https://api.stlouisfed.org/fred"

    def __init__(self, api_key: str | None = None) -> None:
        self.api_key = api_key or get_settings().fred_api_key
        if not self.api_key:
            raise ValueError("FRED_API_KEY is required (set in .env)")

    @http_retry()
    def _get(self, path: str, params: dict) -> dict:
        with limit(self.code, 120, 30):  # 120 req/min, burst 30
            q = {**params, "api_key": self.api_key, "file_type": "json"}
            r = httpx.get(f"{self.base_url}{path}", params=q, timeout=30.0)
            r.raise_for_status()
            return r.json()

    def healthcheck(self) -> bool:
        try:
            data = self._get("/releases", {"limit": 1})
            return "releases" in data
        except Exception:
            return False

    def fetch(self, symbol: str, start, end) -> pl.DataFrame:
        data = self._get(
            "/series/observations",
            {"series_id": symbol, "observation_start": start.isoformat(),
             "observation_end": end.isoformat()},
        )
        rows = [{"date": obs["date"], "value": _parse_value(obs["value"])}
                for obs in data.get("observations", [])]
        return pl.DataFrame(rows, schema={"date": pl.Date, "value": pl.Float64})
\end{lstlisting}

Beachten Sie:

\begin{itemize}
  \item \code{@http\_retry()} -- dekoriert die Methode mit automatischen
        Retries bei Netzwerkfehlern (3 Versuche, 5s/30s/120s Backoff).
  \item \code{with limit(...)} -- Token-Bucket, das 120~req/min
        erlaubt (FREDs Limit) mit Burst~30.
  \item \code{\_parse\_value} -- FRED verwendet \code{"."} für fehlende
        Werte; das wird zu \code{NaN} konvertiert.
\end{itemize}

\section{ECB -- European Central Bank SDMX}

\begin{definition}
\textbf{SDMX (Statistical Data and Metadata Exchange)}
\emph{SDMX} ist ein ISO-Standard für den Austausch statistischer Daten.
Große Institutionen wie EZB, OECD, IWF und Eurostat veröffentlichen
ihre Daten über SDMX-Endpunkte. Das Format hat zwei Häufigkeiten:

\begin{itemize}
  \item \textbf{SDMX-ML} (XML-basiert) -- der offizielle Standard
  \item \textbf{SDMX-JSON} (JSON-basiert) -- moderner, einfacher
        zu parsen
\end{itemize}

Dieses Projekt verwendet die \emph{CSV}-Variante (über
\code{format=csvdata}), weil sie sich am einfachsten in Polars laden
lässt.
\end{definition}

\begin{beispiel}
\textbf{ECB SDMX URL-Struktur}
\begin{lstlisting}[language=shell]
# EXR = Exchange Rates
# D.USD.EUR.SP00.A = Daily, USD, EUR, Spot, Average
https://data-api.ecb.europa.eu/service/data/EXR/D.USD.EUR.SP00.A?startPeriod=2024-01-01&endPeriod=2024-12-31&format=csvdata
\end{lstlisting}
\end{beispiel}

Die ECB-API ist im Gegensatz zu FRED \emph{anonym} -- kein API-Key
nötig. Listing~\ref{lst:ecb-fetch} zeigt den ECB-Vendor.

\begin{lstlisting}[language=python, caption={ECB-Vendor: Parse SDMX CSV},label=lst:ecb-fetch]
class ECBVendor(Vendor):
    code = "ecb"

    def __init__(self, base_url: str | None = None) -> None:
        self.base_url = (base_url or get_settings().ecb_sdmx_base).rstrip("/")

    @http_retry()
    def _get_csv(self, path: str, params: dict) -> str:
        with limit(self.code, 60, 10):
            q = {**params, "format": "csvdata"}
            r = httpx.get(f"{self.base_url}{path}", params=q, timeout=30.0)
            r.raise_for_status()
            return r.text

    def fetch(self, symbol, start, end) -> pl.DataFrame:
        flow, key = self._resolve_symbol(symbol)
        csv_text = self._get_csv(
            f"/data/{flow}/{key}",
            {"startPeriod": start.isoformat(), "endPeriod": end.isoformat()},
        )
        return _parse_ecb_csv(csv_text)
\end{lstlisting}

Die SDMX-Daten haben \emph{zwei verschiedene Datumsformate}:
Tagesdaten verwenden \code{YYYY-MM-DD} (z.\,B.\ \code{2024-06-03}),
Monatsdaten verwenden \code{YYYY-MM} (z.\,B.\ \code{2024-01}).
Der Parser muss beide Formen handhaben.

\begin{lstlisting}[language=python, caption={ECB-CSV-Parser: Beide Datumsformate}]
def _parse_ecb_csv(csv_text: str) -> pl.DataFrame:
    # Disables auto-type-inference (which would mis-parse "2024-01" as broken date)
    df = pl.read_csv(
        csv_text.encode("utf-8"),
        columns=["TIME_PERIOD", "OBS_VALUE"],
        dtypes={"TIME_PERIOD": pl.Utf8, "OBS_VALUE": pl.Float64},
    )
    if df.is_empty():
        return pl.DataFrame(schema={"date": pl.Date, "value": pl.Float64})

    # coalesce: first try full date, fall back to year-month
    df = df.with_columns(
        pl.coalesce(
            pl.col("TIME_PERIOD").str.strptime(pl.Date, format="%Y-%m-%d", strict=False),
            pl.col("TIME_PERIOD").str.strptime(pl.Date, format="%Y-%m", strict=False),
        ).alias("date")
    )
    return df.select(
        pl.col("date"),
        pl.col("OBS_VALUE").cast(pl.Float64, strict=False).alias("value"),
    )
\end{lstlisting}

\section{Yahoo Finance (via yfinance)}

\begin{definition}
\textbf{yfinance}
\emph{yfinance} ist eine Python-Bibliothek, die
\emph{inoffizielle} Yahoo-Finance-API verwendet, um historische
Aktien-, ETF- und Index-Daten abzurufen. Es ist \emph{nicht} offiziell
von Yahoo unterstützt -- die Bibliothek scraped und parst die
Yahoo-Website. Daher kann sie bei Yahoo-Änderungen kaputt gehen,
und die Yahoo \emph{terms of service} (ToS) erlauben
\emph{keine Weiterverteilung} der Daten. Dieses Projekt
verwendet yfinance daher nur für \emph{private} Forschungszwecke.
\end{definition}

\begin{beispiel}
\textbf{yfinance API -- sehr einfach}
\begin{lstlisting}[language=python]
import yfinance as yf
t = yf.Ticker("AAPL")
df = t.history(start="2024-01-01", end="2024-12-31", auto_adjust=False)
# df hat Spalten: Open, High, Low, Close, Volume, Dividends, Stock Splits
# Index: DatetimeIndex mit TZ (z.B. America/New_York)
\end{lstlisting}
\end{beispiel}

\begin{lstlisting}[language=python, caption={Yahoo-Vendor: Single und Bulk Fetch},label=lst:yahoo-fetch]
import yfinance as yf
import pandas as pd

class YahooVendor(Vendor):
    code = "yahoo"

    def fetch(self, symbol, start, end) -> pl.DataFrame:
        with limit(self.code, 30, 5):  # 30 req/min, burst 5
            t = yf.Ticker(symbol)
            pdf = t.history(
                start=start.isoformat(),
                end=(end + timedelta(days=1)).isoformat(),  # yf end is exclusive
                auto_adjust=False,    # we need both close and adj_close
                actions=False,         # splits/dividends go to corporate_actions
            )
        if pdf.empty:
            return pl.DataFrame(schema=_YAHOO_SCHEMA)
        return _yahoo_to_polars(pdf)

    def fetch_many(self, symbols: list[str], start, end) -> dict[str, pl.DataFrame]:
        """Bulk fetch in a single yfinance call -- much faster."""
        with limit(self.code, 30, 5):
            pdf = yf.download(
                tickers=" ".join(symbols),
                start=start.isoformat(),
                end=(end + timedelta(days=1)).isoformat(),
                auto_adjust=False,
                actions=False,
                group_by="ticker",
                threads=True,
                progress=False,
            )
        # ... handle MultiIndex columns for multiple tickers
\end{lstlisting}

Listing~\ref{lst:yahoo-fetch} zeigt den Yahoo-Vendor. Zwei
Dinge sind erwähnenswert:

\begin{enumerate}
  \item \textbf{Bulk-Methode}: \code{fetch\_many} verwendet
        \code{yf.download} mit \code{threads=True} -- dies holt
        \emph{alle} Symbole in einem einzigen API-Call. Für 36~Symbole
        benötigt es nur ~12~Sekunden, statt 36~einzelne Calls.
  \item \textbf{Pandas-zu-Polars-Bridge}: yfinance liefert ein
        \code{pandas.DataFrame} mit TZ-aware-DatetimeIndex und
        nullable Integers (\code{Int64} statt \code{int64}). Wir
        konvertieren explizit zu pl-kompatiblen Typen (\code{int64})
        und parsen die TZ-Datums in Polars.
\end{enumerate}

\section{Rate-Limiting und Retry}

\begin{definition}
\textbf{Exponential Backoff}
\emph{Exponential Backoff} ist eine Retry-Strategie, bei der die
Wartezeit zwischen Versuchen \emph{exponentiell wächst}: 5s, dann
30s, dann 120s. Dies gibt einem überlasteten Server Zeit, sich zu
erholen, ohne dass der Client ihn weiter bombardiert.
\end{definition}

Listing~\ref{lst:retry} zeigt die zentralen Werkzeuge für robustes
HTTP.

\begin{lstlisting}[language=python, caption={Rate-Limiting und Retry-Logik},label=lst:retry]
import httpx
import time
from tenacity import retry, stop_after_attempt, wait_chain, wait_fixed

# Token-Bucket für Rate-Limiting
class TokenBucket:
    def __init__(self, rate_per_min: int, burst: int):
        self.rate = rate_per_min / 60.0  # tokens per second
        self.burst = burst
        self.tokens = float(burst)
        self.last = time.monotonic()

    def take(self, n: int = 1) -> None:
        # Wenn nicht genug Tokens: blockiere bis genug da sind
        now = time.monotonic()
        elapsed = now - self.last
        self.last = now
        self.tokens = min(self.burst, self.tokens + elapsed * self.rate)
        if self.tokens < n:
            time.sleep((n - self.tokens) / self.rate)
            self.tokens = 0
        else:
            self.tokens -= n

# Retry-Decorator mit 5s/30s/120s Backoff
def http_retry():
    return retry(
        reraise=True,
        stop=stop_after_attempt(3),
        wait=wait_chain(wait_fixed(5), wait_fixed(30), wait_fixed(120)),
        retry=retry_if_exception_type((
            httpx.TimeoutException, httpx.NetworkError, httpx.RemoteProtocolError,
        )),
    )
\end{lstlisting}

Jeder Vendor-Adapter verwendet \code{limit(code, rate, burst)} und
\code{@http\_retry()}, um die spezifischen Anforderungen der Quelle
zu erfüllen. Tabelle~\ref{tab:rate-limits} zeigt die Limits pro
Vendor.

\begin{table}[ht]
\centering
\small
\begin{tabular}{lrrl}
  \toprule
  \textbf{Vendor} & \textbf{Rate} & \textbf{Burst} & \textbf{Retry-Backoff} \\
  \midrule
  FRED   & 120 req/min & 30 & 5s/30s/120s \\
  ECB    & 60 req/min  & 10 & 5s/30s/120s \\
  Yahoo  & 30 req/min  & 5  & 5s/30s/120s \\
  \bottomrule
\end{tabular}
\caption{Rate-Limits und Burst pro Vendor. Diese sind in
  \code{src/deephold\_db/vendors/<name>.py} konfiguriert.}
\label{tab:rate-limits}
\end{table}

\section{Was als Nächstes?}

Die Vendor-Adapter sind die Datenquellen; die ETL-Pipeline
(Kapitel~\ref{ch:etl}) zeigt, wie diese Daten in unsere
Datenbank gelangen und für Analysen verfügbar gemacht werden.
